\documentclass[../DoAn.tex]{subfiles}
\begin{document}
\section{Kết luận}
Trong khi tiến hành thiết kế, xây dựng và triển khai hệ thống ecommerce cho cả mobile và web admin này, vì để thực hiện các chức năng đáp ứng các nhu cầu của hệ thống, bản thân em cũng đã được tiếp xúc với nhiều loại công nghệ mới mà em chưa bao giờ được biết trước đây. Nhờ vậy, mà em đã xây dựng thành công một hệ thống có những đặc điểm sau:
\begin{itemize}
    \item Xây dựng thành công ứng dụng flutter mobile thương mại điện tử với giao diện đẹp mắt, thân thiện, đầy đủ các tính năng cơ bản của một hệ thống thương mại điện tử, ngoài ra đã phát triển được một số tính năng, đóng góp nổi bật, có sự khác biệt với các sản phẩm trên thị trường như tính năng chatbot tự động hóa quy trình mua sắm, gợi ý cá nhân hóa theo sở thích, lịch sử người dùng, tính năng săn sale theo nhóm, hệ thống nhiệm vụ đầy thử thách. 
    \item Deploy được một số backend API như backend cho train tự động mô hình gợi ý, API cho chatbot \cite{RasaFramework} gọi đến lên \cite{GoogleCloudRun}, có lên lịch tự động chạy các API.
     \item Sử dụng các dịch vụ của nền tảng Firebase từ Firestore để lưu trữ database~\cite{FirebaseFirestore}, 
  Firebase Authentication để xác thực đăng nhập bằng Google~\cite{FirebaseAuth}, 
  Firebase Storage để lưu trữ ảnh, video~\cite{FirebaseStorage}, 
  Firebase Cloud Messaging để gửi thông báo realtime~\cite{FirebaseCloudMessaging}, 
  và Firebase App Distribution để người dùng tải ứng dụng và trải nghiệm~\cite{FirebaseAppDistribution}. 
  Đây đều là những kiến thức và trải nghiệm mới đầy thú vị.
  \item Xây dựng được trang web dành cho quản trị viên với giao diện đẹp mắt sử dụng \cite{ReactJS}, backend sử dụng \cite{SpringBoot} 
  \item Khả năng tự nghiên cứu và khai thác các mô hình AI hiện đại, đặc biệt là mô hình gợi ý LightFM~\cite{LightFM}, thông qua nền tảng Kaggle~\cite{Kaggle} và ngôn ngữ lập trình Python~\cite{Python}, từ đó từng bước xây dựng quy trình pipeline để xử lý dữ liệu, huấn luyện mô hình và đánh giá hiệu quả gợi ý. Việc áp dụng thư viện \cite{LightFM} không chỉ giúp tăng độ chính xác trong gợi ý mà còn mở rộng kiến thức về hybrid recommendation systems kết hợp collaborative filtering và content-based filtering.

\end{itemize}

Tuy nhiên, song song với những điểm mạnh ấy, ứng dụng vẫn còn một số những khuyết điểm như:

\begin{itemize}
    \item Một số request, truy vấn database còn khá chậm ảnh hưởng đến trải nghiệm người dùng, cần tối ưu nhất có thể để mọi response của hệ thống chỉ phản hồi tối đa 1-2s.
    \item Ứng dụng hiện chỉ mới được xây dựng dành riêng cho thiết bị điện thoại thông thường, giao diện chưa tối ưu và thích ứng tốt trên các phiên bản màn hình lớn như tablet.
    \item Chatbot hiện tại vẫn hiểu được khá ít các mẫu câu tự nhiên của người dùng, cần huấn luyện thêm để phản hồi tự nhiên, chính xác hơn với nhiều câu nhập khác nhau của người dùng hơn
    \item Cần thu thập nhiều data hơn, so sánh thêm nhiều mô hình gợi ý khác ngoài\cite{LightFM} để càng ngày càng cải thiện tính năng gợi ý tốt hơn nữa
\end{itemize}

% Sinh viên so sánh kết quả nghiên cứu hoặc sản phẩm của mình với các nghiên cứu hoặc sản phẩm tương tự.

% Sinh viên phân tích trong suốt quá trình thực hiện ĐATN, mình đã làm được gì, chưa làm được gì, các đóng góp nổi bật là gì, và tổng hợp những bài học kinh nghiệm rút ra nếu có.

\section{Hướng phát triển}

\subsection{Tính năng Livestream kết hợp Đấu giá ngược}
Trong tương lai, em mong muốn tích hợp cơ chế \textbf{đấu giá ngược} vào tính năng livestream để nâng cao mức độ tương tác giữa người bán và người mua. Cụ thể, trong quá trình livestream, người bán có thể kích hoạt chế độ giảm giá dần theo thời gian cho một sản phẩm cụ thể. Người mua có thể theo dõi giá giảm theo thời gian thực và quyết định “chốt đơn” tại thời điểm phù hợp. Càng chờ lâu, giá càng rẻ – nhưng nguy cơ sản phẩm bị người khác mua trước cũng cao hơn. Đây là một cơ chế giải trí, kích thích FOMO (sợ bỏ lỡ) và tạo trải nghiệm mua sắm thú vị, khác biệt hoàn toàn so với hình thức flash sale truyền thống. Việc triển khai đấu giá ngược cần kết hợp giữa cập nhật real-time, UI động và xử lý đồng thời nhiều phiên livestream khác nhau.

\subsection{Tính năng Trợ lý ảo gợi ý phối đồ (AI Stylist)}
Để nâng cao trải nghiệm cá nhân hóa trong lĩnh vực thời trang, em hướng tới phát triển một \textbf{trợ lý ảo gợi ý phối đồ} sử dụng AI. Tính năng này sẽ học từ lịch sử mua hàng, hành vi xem sản phẩm và xu hướng thời trang để đưa ra gợi ý phối đồ phù hợp theo mùa, sự kiện hoặc phong cách riêng của người dùng. Kết hợp với công nghệ thực tế tăng cường (AR), người dùng còn có thể thử đồ trực tiếp trên camera thông qua “phòng thử đồ ảo”. Đây là bước đột phá giúp cá nhân hóa trải nghiệm người dùng ở mức cao, đồng thời góp phần giảm tỷ lệ hoàn trả sản phẩm vì không vừa ý. Việc xây dựng hệ thống cần tích hợp recommendation system, xử lý hình ảnh thời gian thực và UI AR hiện đại.

\subsection{Tính năng Trợ lý mua sắm thông minh tích hợp AI}
Bên cạnh các chức năng cơ bản, một hướng đi mới đầy tiềm năng là xây dựng \textbf{trợ lý mua sắm thông minh} (Smart Shopping Assistant) tích hợp AI để đồng hành cùng người dùng trong toàn bộ hành trình mua sắm. Trợ lý này có thể:
\begin{itemize}
    \item Phân tích thói quen mua sắm và đưa ra gợi ý phù hợp theo nhu cầu.
    \item Gợi ý thời điểm mua tốt nhất dựa trên lịch sử giá sản phẩm.
    \item Nhắc người dùng nếu sản phẩm yêu thích sắp hết hàng hoặc đang có khuyến mãi.
    \item Hỗ trợ người dùng tìm sản phẩm thông qua giọng nói hoặc hình ảnh.
\end{itemize}
Việc triển khai tính năng này cần kết hợp nhiều thành phần công nghệ như NLP (Natural Language Processing) để xử lý ngôn ngữ tự nhiên, computer vision để nhận diện hình ảnh sản phẩm, và mô hình học máy để dự đoán hành vi người dùng. Đây sẽ là yếu tố then chốt nâng tầm trải nghiệm cá nhân hóa và đưa nền tảng tới gần hơn với xu thế thương mại điện tử thông minh (Smart e-Commerce).

\end{document}